\documentclass[11pt]{article}

\usepackage[margin=0.82in]{geometry}
\usepackage{amsmath,amssymb,booktabs,array,longtable,enumitem,microtype,xcolor,fancyhdr,lastpage}
\usepackage[T1]{fontenc}
\usepackage{lmodern}
\usepackage{titlesec}
\usepackage{parskip}
\usepackage{caption}
\usepackage{hyperref}

\hypersetup{colorlinks=true,linkcolor=black,urlcolor=blue,citecolor=black}

\definecolor{darkblue}{RGB}{26,54,93}

\titleformat{\section}{\large\bfseries\color{darkblue}}{\thesection}{0.6em}{}
\titleformat{\subsection}{\normalsize\bfseries}{\thesubsection}{0.6em}{}
\titleformat{\subsubsection}{\normalsize\bfseries}{\thesubsubsection}{0.6em}{}

\setlist[itemize]{leftmargin=1.3em,itemsep=0.2em,topsep=0.25em}
\setlist[enumerate]{leftmargin=1.5em,itemsep=0.2em,topsep=0.25em}
\setlength{\tabcolsep}{3pt}
\setlength{\LTleft}{0pt}
\setlength{\LTright}{0pt}

\pagestyle{fancy}
\fancyhf{}
\lhead{Aggression Monitor V1}
\rhead{Implemented measurement}
\cfoot{\thepage\ / \pageref{LastPage}}
\setlength{\headheight}{14pt}

\newcommand{\code}[1]{\texttt{\detokenize{#1}}}
\newcommand{\eps}{\varepsilon}

\title{\vspace{-1.5cm}\textbf{Aggression Monitor V1}\\[0.25em]
\large What the live page measures, how each chart is drawn, and how the
dominance banner is classified}
\author{}
\date{Implementation as of 22 August 2026}

\begin{document}
\maketitle
\thispagestyle{fancy}
\vspace{-1.2em}

\begin{abstract}
This note describes only Aggression Monitor V1
(\code{scripts/aggression_dashboard/}, port 5050,
\url{https://aasthajain.tail174727.ts.net/}).
The page is a live measurement of curve aggressive flow against a chosen
Treasury-futures price.  It is not a trading system.  It does not emit
long or short, RANGE or TREND, alerts, or episode scores.

There \emph{are} two fitted lines labeled fair value.  They are not
multi-feature forecasts of future price.  Each is a contemporaneous
ordinary-least-squares fit of price (or same-day price change) on
cumulative TotalAggress.  The residual of that fit is then reused by the
flow-dominance banner as ``implied passive.''  Implied passive is
\(\Delta e\), the change in that residual.  It is not resting-order
identity.
\end{abstract}

\tableofcontents
\newpage

\section{Goal}

V1 answers a measurement question: how much aggressive buying and selling
is occurring across the Treasury curve, and how does that stack up
against the chosen contract's price path in the same minutes?

The five charts and the top banner are different views of that same
construction:

\begin{itemize}
    \item Panels 1--3 show the raw flow and price together.
    \item Panels 4--5 draw a fitted line of price on TotalAggress and mark
    large residuals.
    \item The banner's headline is one of AGGRESS DOMINANT, IMPLIED
    PASSIVE DOMINANT, BALANCED, or QUIET (or WARMING / UNAVAILABLE).
    A second line reports z-scores, whether each leg is increasing /
    slowing / steady, and what the headline just shifted from.
\end{itemize}


\section{Page, controls, and data}

\subsection{What the user selects}

Default live view: price symbol TY, aggression symbols FV, TU, TY, and US
all on, DV01 adjust on, lookback 2 hours, start time the current trading-day
open (18:00 ET),
auto-refresh on, live mode.

\begin{center}
\begin{tabular}{@{}ll@{}}
\toprule
Control & Effect \\
\midrule
Price symbol & Which contract's price is drawn on the right-hand axes
(FV, TU, TY, US; default TY). \\
Aggression chips & Which of those four roots contribute flow
(default: all four). \\
DV01 Adjust & If on, each contract's size is multiplied by a fixed
dollar DV01 before summing. \\
Lookback & How much of the built series the charts display:
2, 6, 12, or 24 hours. \\
Live vs historical & Live merges the AccessAPI buffer with a session
backfill.  Historical pulls a fixed window. \\
\bottomrule
\end{tabular}
\end{center}

The default start is the current trading-day open, 18:00 ET.

\subsection{Sources}

Live trades come from AccessAPI
(\code{ws://100.79.30.21:8001/ws/live/stream}) into an 18-hour in-memory
buffer.  If that feed is stale for 30 seconds, a Databento live TBBO
backup can start; it stops after 60 seconds of primary recovery.
Historical and backfill trades are Databento parent-symbol trade
schema pulls for ZF, ZT, ZN, and ZB.  Price is the last PostgreSQL OHLCV
print in each one-minute bucket, forward-filled, with trade price as
fallback.

A background thread refreshes session history from the trading-day open every
240 seconds, overlapping the last 15 minutes.  Trade-cache parquet older
than 30 days is pruned.  Front-month selection, per UTC calendar date, keeps
the \code{instrument_id} with the highest traded size that day.

Trades with non-positive size, or price outside \([80,150]\), are dropped.

\subsection{One-minute grid}

All series are resampled to one-minute buckets
(\code{DEFAULT_BUCKET=1min}).  The latest bucket can still be forming
except in the dominance banner, which drops it.

\section{How a trade becomes flow}

\subsection{Aggressor sign}

Each trade has a source side.  The mapping is

\begin{equation}
\operatorname{sign}
=
\begin{cases}
+1, & \text{side }\in\{B,\mathrm{BID}\},\\
-1, & \text{side }\in\{A,\mathrm{ASK}\},\\
0, & \text{otherwise}.
\end{cases}
\end{equation}

\(+1\) is an aggressive buy (lifted the offer).  \(-1\) is an aggressive
sell (hit the bid).  Unknown sides contribute zero flow and are counted
in metadata only.  There is no quote-rule fallback.



\subsection{Bucket sums}

Over one-minute bucket \(t\), summing selected symbols,

\begin{align}
M_t^b
&=
\sum_{i\in t}\mathbf{1}\{\operatorname{sign}_i=+1\}\,\mathrm{size}_i\,\mathrm{DV01}_{c(i)},\\
M_t^s
&=
\sum_{i\in t}\mathbf{1}\{\operatorname{sign}_i=-1\}\,\mathrm{size}_i\,\mathrm{DV01}_{c(i)},\\
\Delta A_t
&=
M_t^b-M_t^s.
\end{align}

\(M_t^b\) is Buy Aggress.  \(M_t^s\) is Sell Aggress.  \(\Delta A_t\) is
the signed one-minute increment.

\subsection{TotalAggress}

TotalAggress is the running sum of the one-minute increments
\(\Delta A_t\), starting at the trading-day open (18:00 ET):

\begin{equation}
A_t
=
\sum_{u=t_{\mathrm{open}}}^{t}\Delta A_u.
\end{equation}

Lookback only shortens the time axis.  It does not restart the sum, so
the last two hours still show the session-to-date level.


\subsection{BuySell Ratio}

The raw imbalance is a signed percentage, not buy volume divided by sell
volume:

\begin{equation}
\mathrm{BuySellRaw}_t
=
100\cdot
\frac{\Delta A_t}{M_t^b+M_t^s}
\quad\text{when }M_t^b+M_t^s\ne 0.
\end{equation}

The plotted series is an EMA of that raw ratio with span 10 and
\code{ignore_na=True}, \code{min_periods=1}.  Empty-volume minutes stay
missing in the raw series and do not force the EMA to zero.

\section{The five charts}

Times on the \(x\)-axis are Eastern, timezone-naive ISO strings.
Lookback trims the time axis only.  The TotalAggress level stays
session-to-date.

\subsection{Panel 1. TotalAggress vs Price/Yield}

Two lines, dual axis.  Left: session TotalAggress (blue).  Right: the
selected price (red).  Price is last-in-bucket, forward-filled onto the
flow grid by \code{align_flow_and_price}.

The dual-axis ranges are scaled so the two series occupy the same
vertical box, with a 5\% margin.  That visual \(\alpha,\beta\) is
\emph{not} the fair-value OLS.  It is only a Plotly range mapping.

Read: when the blue cumulative flow and the red price travel together,
aggressive flow and price are moving in the same direction in the
displayed window.  When they diverge, price is moving for reasons other
than this cumulative aggression.

\subsection{Panel 2. BuySell Ratio vs Price/Yield}

Left: the EMA BuySell Ratio (green).  Right: the same price (orange).
Ratio is forward-filled onto the price grid.

Read: a positive ratio means that minute's signed flow was buy-heavy; a
negative ratio is sell-heavy.  Because of the EMA, a single burst is
smoothed across nearby minutes.

\subsection{Panel 3. Buy Aggress, Sell Aggress, TotalAggress, and price}

Four traces on three axes:

\begin{itemize}
    \item green bars: Buy Aggress \(M_t^b\);
    \item red bars: Sell Aggress, plotted as \(-|M_t^s|\) so sells hang
    below zero;
    \item black line: TotalAggress on a third axis;
    \item brown line: price on the second axis.
\end{itemize}

Read: the bars are the one-minute activity; the black line is the
running sum; the brown line is price.  This is the only panel that shows
buys and sells separately.

\subsection{Panel 4. Intraday Fair Value (fit once per session)}

Same axis: actual price (solid blue) and fitted price (dashed red).
Grey markers sit on actual price where \(|z|\ge 2\).

This is the first of the two ``fair value'' charts.  \textbf{It is a
contemporaneous fit, not a forecast}.  After at least 30 aligned minutes
with a non-constant session TotalAggress \(A_t\),

\begin{equation}
P_t=\alpha+\beta A_t+e_t.
\end{equation}

The dashed line is \(\widehat{P}_t=\alpha+\beta A_t\).
The residual is \(e_t=P_t-\widehat{P}_t\).  The displayed z-score is

\begin{equation}
z_t^{\mathrm{V1,live}}=e_t/s_e,
\end{equation}

where \(s_e\) is the sample standard deviation of residuals in the
currently aligned frame.  The residual mean is not subtracted.
A point is marked large when \(|z_t^{\mathrm{V1,live}}|\ge 2\).

The first successful fit is cached under the current trading day,
price symbol, aggression-symbol tuple, DV01 flag, and bucket size.
Later refreshes reuse \(\alpha,\beta\) until that key changes.

Read: the dashed line is ``where price would sit if it were a linear
function of this session's cumulative aggression.''  When solid price is
above the dash, the residual is positive: price is rich relative to that
line.  That is a description of the current sample, not a prediction
that price will revert.

\subsection{Panel 5. 10-Business-Day Fair Value}

Same drawing language as panel 4, with red large-residual markers.
The source window is the last 10 business days from 18:00 ET on
the first of those dates through now.  That archive is refreshed at most
once per trading day.

Each trading day first subtracts that session's opening price
\(P_{\mathrm{open},s}\).  One pooled OLS is then fit:

\begin{equation}
P_t-P_{\mathrm{open},s}
=
\alpha_{10}+\beta_{10}A_t+e_t^{(10)},
\qquad
\widehat{P}^{(10)}_t
=
P_{\mathrm{open},s}+\alpha_{10}+\beta_{10}A_t.
\end{equation}

Read: same idea as panel 4, but \(\beta_{10}\) is estimated across ten
trading days of session-reset flow versus session price change.  It
is still a contemporaneous line, not a walk-forward forecast.

\section{The flow-dominance banner}

The live banner has three lines.  The first line is the state.  The
second is a status sentence.  The third is a fixed disclaimer plus an
as-of time.

The headline states, exactly as drawn, are

\begin{center}
\textsc{AGGRESS DOMINANT}, \quad
\textsc{IMPLIED PASSIVE DOMINANT}, \quad
\textsc{BALANCED}, \quad
\textsc{QUIET}.
\end{center}

Implied means leftover price change after the fitted aggression piece
(\(\Delta e\)), not observed resting-order flow.

Before the classification is ready the headline is WARMING or
UNAVAILABLE.  Internally those four ready states are
\code{AGGRESS_DOM}, \code{PASSIVE_DOM}, \code{NEITHER}, and
\code{QUIET}.

When status is READY, the second line is built as

\begin{center}
\texttt{shifted from <prior state> <n>m ago} \quad\textrm{or}\quad
\texttt{holding <n>m}\\
\texttt{· aggress increasing|slowing|steady (z +x.xx)}\\
\texttt{· implied passive increasing|slowing|steady (z +x.xx)}.
\end{center}

The two z values are \(z_A\) and \(z_P\) from the 15-row contribution
windows, not panel 4's residual z-score.  If the banner is not READY,
that second line is a reason such as waiting for completed one-minute
data.


The banner uses the \emph{intraday} fit from panel 4, on completed
one-minute bars only.  The forming bucket is excluded.  If that model is
missing, the headline is UNAVAILABLE.

\subsection{Contributions}

On each trading day, one-minute contributions are

\begin{align}
\Delta\mathrm{Aggress}_t
&=
\beta\cdot\Delta A_t,\\
\Delta\mathrm{Passive}_t
&=
\Delta e_t,\\
\Delta P_t
&=
\Delta P_t.
\end{align}

\(\Delta\mathrm{Aggress}\) is the piece of the price change attributed
to the fitted slope times the change in TotalAggress.
\(\Delta\mathrm{Passive}\) is the change in the residual.  The page
calls that second piece implied passive.  It is not observed book
additions or cancellations.

Each series is then summed over a trailing window of exactly 15
completed rows on that trading day (not necessarily 15 wall-clock minutes
if rows are missing):

\begin{equation}
S^A_t=\sum_{j=0}^{14}\Delta\mathrm{Aggress}_{t-j},
\qquad
S^P_t=\sum_{j=0}^{14}\Delta\mathrm{Passive}_{t-j},
\qquad
\Delta P_{15,t}=\sum_{j=0}^{14}\Delta P_{t-j}.
\end{equation}

Those window sums are standardized against earlier valid 15-row windows
on the same trading day.  Let \(\mu_{t^-}\) and \(\sigma_{t^-}\) be the
expanding mean and sample standard deviation of earlier windows, needing
at least 20 prior valid windows.  Then

\begin{equation}
z_A=(S^A_t-\mu^A_{t^-})/\sigma^A_{t^-},
\qquad
z_P=(S^P_t-\mu^P_{t^-})/\sigma^P_{t^-}.
\end{equation}

These z-scores are not \(z^{\mathrm{V1,live}}\) from panel 4.  Until 20
prior windows exist, the state is WARMING.

\subsection{State rules}

Fixed parameters: window 15 rows, dominance ratio 1.5, quiet floor 0.5.

\begin{itemize}
    \item \textsc{quiet} if \(|z_A|<0.5\) and \(|z_P|<0.5\).
    \item \textsc{aggress dominant} if \(\operatorname{sign}(z_A)=\operatorname{sign}(\Delta P_{15})\ne 0\)
    and \(|z_A|\ge 1.5\,|z_P|\).
    \item \textsc{implied passive dominant} if \(|z_P|\ge 1.5\,|z_A|\) and
    either \(z_P\) and \(z_A\) have opposite nonzero signs, or
    \(\operatorname{sign}(z_P)=\operatorname{sign}(\Delta P_{15})\ne 0\).
    \item \textsc{balanced} otherwise (internally \code{NEITHER}).
\end{itemize}

Displayed names: AGGRESS DOMINANT, IMPLIED PASSIVE DOMINANT, BALANCED,
QUIET, plus WARMING and UNAVAILABLE.

In words:

\begin{itemize}
    \item Quiet: both legs are small relative to today so far.
    \item Aggress dominant: the aggression-attributed piece is the
    stronger one and it has the same sign as the 15-row price change.
    \item Implied passive dominant: the residual piece is the stronger
    one, and it is either fighting the aggression piece or agreeing with
    the 15-row price change.
    \item Balanced: neither quiet nor a 1.5-to-1 winner under those
    sign rules.
\end{itemize}

\subsection{Dynamics line}

The detail line compares each leg's absolute z with its value five
minutes earlier (or the latest valid z at or before that timestamp).
A change larger than \(0.15\) in absolute z is increasing; a change
smaller than \(-0.15\) is slowing; otherwise steady.  Missing history
is WARMING.

It also reports whether the classified state just shifted and how many
completed minutes the current state has lasted.

The banner does not persist episodes, does not score outcomes, and does
not write a live store.

\section{Parameter register}

\begin{longtable}{@{}p{0.34\linewidth}p{0.20\linewidth}p{0.38\linewidth}@{}}
\toprule
Name & Value & Role \\
\midrule
\endfirsthead
\toprule
Name & Value & Role \\
\midrule
\endhead
Price default & TY & Right-hand price series. \\
Aggression default & FV, TU, TY, US & Flow contributors. \\
Bucket & 1 minute & All series. \\
TotalAggress & sum from 18:00 ET open & Same path on charts and both fits. \\
DV01 FV/TU/TY/US & 42 / 36 / 85 / 125 & Dollars per bp; default on. \\
BuySell EMA span & 10 & Displayed ratio smoothing. \\
Default start & 18:00 ET trading-day open & Current Globex session. \\
Default lookback & 2 hours & Chart trim. \\
Live buffer retention & 18 hours & AccessAPI memory. \\
Backup stale / recover & 30s / 60s & Databento failover. \\
Backfill period / overlap & 240s / 900s & Session history refresh. \\
Trade-cache retention & 30 days & Historical parquet. \\
Trade price filter & \([80,150]\) & Dropped otherwise. \\
Front-month rule & max size by UTC date & One \code{instrument_id} per root. \\
Intraday min points & 30 & OLS will not fit before this. \\
Intraday large residual & \(|z|\ge 2\) & Grey markers on panel 4. \\
Intraday cache & once per trading day & Reuse \(\alpha,\beta\) until the session rolls. \\
10-day window & 10 business days from 18:00 ET & Panel 5 source. \\
10-day refresh & once per trading day & Source refetch hourly if session rolled. \\
10-day large residual & \(|z|\ge 2\) & Red markers on panel 5. \\
Banner window & 15 completed rows & Contribution sums. \\
Banner dominance ratio & 1.5 & Winner threshold. \\
Banner quiet floor & 0.5 & Both \(|z|\) below this. \\
Banner z warmup & 20 prior windows & Else WARMING. \\
Banner dynamics lookback & 5 minutes & Increasing / slowing / steady. \\
Banner dynamics deadband & 0.15 & Ignore smaller \(|z|\) changes. \\
Banner completed bars & forming bar dropped & Unlike panels 1--5. \\
\bottomrule
\end{longtable}

\section{What we are watching}

The page is for one question: when the chosen contract moves, is that
move in the aggressive curve flow, or is price doing something the
prints do not explain?

The page does not mark RANGE or TREND.  The types below are only ways
to read the same charts and banner.

\subsection{Flow-explained move}

Price and session TotalAggress travel together (panel 1), and the banner
is AGGRESS DOMINANT with aggress increasing.  The dashed fair-value line
stays close to solid price.  The working idea: the move is in the
aggressive flow.

\subsection{Aggress-led breakout}

Price leaves a recent range while TotalAggress has been rising or
falling for many minutes, BuySell stays one-sided, and the banner stays
AGGRESS DOMINANT / increasing.  The working idea: a breakout that is
being paid for in aggressive size, not a thin spike.

\subsection{Implied-passive range}

Price chops or stalls while the banner is IMPLIED PASSIVE DOMINANT or
BALANCED, and the residual \(e\) is swinging rather than trending.
The working idea: aggressive flow is there, but leftover price change
is fighting it, so the path stays contained.

\subsection{Unexplained extension}

Price keeps going while TotalAggress flattens or turns, the residual
grows (\(|z|\ge 2\) on panel 4), and implied passive is increasing.
The working idea: the continuation is not in the curve prints.
Something else is carrying price.

\subsection{Absorption / failed push}

Aggress is large and one-sided, but price barely moves or moves the
other way; implied passive is the stronger leg and opposite in sign.
The working idea: aggressive flow is being taken the other way.  That
is often a range or reversal setup, not a breakout.

\subsection{Exhaustion after a paid run}

A long stretch of AGGRESS DOMINANT / increasing, then aggress goes
slowing or QUIET while price is still extended from the dashed line.
The working idea: the flow that built the move has faded; the next
state is more likely a pause or give-back than a fresh breakout.

\subsection{Quiet then ignition}

The banner is QUIET, then flips to AGGRESS DOMINANT as \(\Delta A\)
jumps and price starts to run.  The working idea: a new impulse, not
the middle of an old one.  Early in the session this can just be
warmup ending.

\subsection{Session vs 10-day disagreement}

Panel 4 (this session) and panel 5 (ten trading days) tell different
residual stories: one marks large \(z\), the other does not.  The
working idea: this session's price-to-flow slope is unlike the recent
sample.  Treat rich or cheap versus flow as session-specific until they
agree.

\subsection{Curve vs chosen price}

TotalAggress includes FV, TU, TY, and US, but the right-hand price is
only one contract.  If TY is flat while the blue sum is moving, the
flow may be in the wings.  The working idea: a TY unexplained residual
can just be other contracts dominating the sum.

These are reading habits, not validated states.  The banner never
becomes a long or short.

\end{document}
